带电粒子散射会改变其长程 Coulomb 场，因此任意低频电磁辐射与硬散射态不可完全分离。探测器只能以有限能量分辨率区分辐射：若未分辨光子的总能量低于阈值 $\Delta E_{\rm det}$，实验可观测量必须把这些末态与无辐射硬末态共同计入。红外有限性因此是包容测量定义与 QED 幺正性的共同结果，而不是单张图的性质。

带电粒子的速度发生改变时，其长程电磁场必须重新调整，这一调整包含频率可以任意低的辐射。严格要求末态“没有任何光子”会把实验上不可分辨的低频辐射排除在外，从而产生红外发散。有限的 $\Delta E_{\rm det}$ 把所有不可分辨软光子态组成同一个实验结果。


设硬过程的约化振幅为 $\mathcal M_0$，外部带电腿编号 $i=1,\ldots,N$。第 $i$ 条腿的电荷写成 $Q_i e$，其中 $Q_i$ 是带符号的无量纲电荷数；出射腿取 $\sigma_i^{\rm io}=+1$，入射腿取 $\sigma_i^{\rm io}=-1$。再发射一条物理四动量为
$k^\mu=(\hbar\omega/c,\hbar\vb{k})$、横向偏振为 $\varepsilon_\lambda^\mu(k)$ 的软光子。令 $p_{\rm hard}$ 表示硬过程的典型物理动量尺度。当
\[
 \epsilon_{\rm soft}\equiv
 \max_i\frac{|p_i\!\cdot k|}{p_{\rm hard}^2}\ll1
\]
时，内部离壳传播子对 $k\to0$ 保持有限，只有外部在壳带电腿产生 $1/(p_i\!\cdot k)$ 极点。以一条出射费米子腿为例，软光子附着在该腿上的精确外线因子可写成
\begin{equation}
 {
 \mathcal L_{i,\lambda}^{\rm out}(k)
 =g_{\rm em}Q_i\,
 \bar u(p_i)\slashed\varepsilon_\lambda^*(k)
 \frac{\slashed p_i+\slashed k+m_ic}
 {(p_i+k)^2-m_i^2c^2+\ii0^+}.}
 \label{eq:soft-external-leg-kernel-dp68}
\end{equation}
它仍保留完整的 $k$ 依赖；软近似只在下一步展开这个因子。\begin{theorem}[带电 Dirac 外腿的领先软光子定理]
考虑树级硬散射，所有带电 Dirac 外腿均有非零质量，硬振幅光滑且电荷守恒，内部线不处于额外在壳或阈值奇点。
取 $k=\tau q$、$\tau\to0^+$，$q$ 为固定正能非零零模四动量，偏振有光滑横向极限；硬动量按守恒约束平滑延拓。
则以下领先软因子化成立。回路修正及严格红外极限须重新核查解析性和余项，本证明不自动给出它们的 $\order(\omega^0)$ 估计。
领先软振幅因此因子化为
\begin{equation}
 {
 \mathcal M_{1\gamma}^{\rm soft}(k,\lambda)
 =S_\lambda(k)\,\mathcal M_0+\order(\omega^0),
 \qquad
 S_\lambda(k)
 =g_{\rm em}\sum_{i=1}^{N}
 \sigma_i^{\rm io}Q_i
 \frac{p_i\!\cdot\varepsilon_\lambda^*(k)}{p_i\!\cdot k}.}
 \label{eq:weinberg-soft-photon-factor-dp10}
\end{equation}
这里 $g_{\rm em}=\sqrt{4\pi\alpha}$ 无量纲。由于增加一条外部光子会改变约化振幅的质量维数，$S_\lambda$ 本身带相应的逆动量量纲；与光子相空间组合后的辐射概率才是无量纲量。硬过程的全部内部动力学只保留在 $\mathcal M_0$ 中。电荷守恒
$\sum_i\sigma_i^{\rm io}Q_i=0$ 使 $\varepsilon^\mu\to\varepsilon^\mu+\zeta k^\mu$ 不改变\eqnrefs{eq:weinberg-soft-photon-factor-dp10}，因此领先软因子自身满足规范不变性。
\end{theorem}
\begin{proof}
令 $k=\tau q$，正能有质量外腿满足 $p_i\cdot q>0$。
硬子振幅的平滑延拓为 $\mathcal M_0+\order(\tau)$，所设内部非奇异传播只贡献 $\order(1)$。
因此领先奇异性仅需从以下外腿计算提取。
由 $p_i^2=m_i^2c^2$ 与 $k^2=0$，分母为

\begin{equation*}
 (p_i+k)^2-m_i^2c^2=2p_i\!\cdot k.
\end{equation*}
分子中的质量壳奇异部分由 Clifford 代数换序分离：
\begin{align*}
 \bar u_i\slashed\varepsilon^*(\slashed p_i+\slashed k+m_ic)
 &=\bar u_i\slashed\varepsilon^*\slashed p_i
   +\bar u_i\slashed\varepsilon^*\slashed k
   +m_ic\,\bar u_i\slashed\varepsilon^*,\\
 \bar u_i\slashed\varepsilon^*\slashed p_i
 &=2p_i\!\cdot\varepsilon^*\,\bar u_i
   -\bar u_i\slashed p_i\slashed\varepsilon^*.
\end{align*}
再用在壳 Dirac 方程 $\bar u_i\slashed p_i=m_ic\,\bar u_i$，质量项相消，留下
\begin{equation*}
 \bar u_i\slashed\varepsilon^*(\slashed p_i+\slashed k+m_ic)
 =2p_i\!\cdot\varepsilon^*\,\bar u_i
 +\bar u_i\slashed\varepsilon^*\slashed k.
\end{equation*}
最后一项为 $\order(\omega)$，除以 $p_i\!\cdot k=\order(\omega)$ 后只贡献 $\order(\omega^0)$；领先项因此为
\begin{equation*}
 g_{\rm em}Q_i
 \frac{p_i\!\cdot\varepsilon_\lambda^*}{p_i\!\cdot k}\,\bar u_i.
\end{equation*}
入射外腿的相邻传播子动量为 $p_i-k$，故
\begin{equation*}
 (p_i-k)^2-m_i^2c^2=-2p_i\!\cdot k,
\end{equation*}
从而相对多一个负号。对反费米子外腿，顶角的物理电荷本身已经带有与相应费米子相反的符号；其外线旋量满足
\[
 (\slashed p_i+m_ic)v(p_i)=0,
 \qquad
 \bar v(p_i)(\slashed p_i+m_ic)=0.
\]
将这两个在壳关系代入与上面相同的分子换序，质量项仍然抵消，领先分子仍化为 $2p_i\!\cdot\varepsilon^*$。因此外腿方向只通过传播子分母的符号进入，粒子种类的电荷符号则由 $Q_i$ 本身进入。定义 $\sigma_i^{\rm io}=+1$ 表示出射外腿、$\sigma_i^{\rm io}=-1$ 表示入射外腿后，每条带电外腿统一贡献
\[
 g_{\rm em}\,\sigma_i^{\rm io}Q_i
 \frac{p_i\!\cdot\varepsilon_\lambda^*}{p_i\!\cdot k}.
\]
对全部外腿求和即得\eqnrefs{eq:weinberg-soft-photon-factor-dp10}。

规范位移 $\varepsilon_\lambda^\mu\mapsto\varepsilon_\lambda^\mu+\zeta k^\mu$ 使软因子改变
\begin{equation*}
 \delta S=g_{\rm em}\zeta^*\sum_i\sigma_i^{\rm io}Q_i.
\end{equation*}
硬过程的电荷守恒给 $\sum_i\sigma_i^{\rm io}Q_i=0$，所以 $\delta S=0$。这同时检验了领先软因子的规范不变性。
\end{proof}

实软辐射的概率由\eqnrefs{eq:weinberg-soft-photon-factor-dp10} 与一光子 Lorentz 不变相空间共同决定。记 $\vb{n}$ 为光子传播方向，$\boldsymbol\beta_i=c\vb{p}_i/E_i$，并定义相对于 $\vb{n}$ 的横向速度
$\boldsymbol\beta_i^\perp=\boldsymbol\beta_i-\vb{n}(\vb{n}\!\cdot\boldsymbol\beta_i)$。用红外能量调节参数 $\lambda_{\rm IR}>0$ 暂时截断 $\hbar\omega\to0$ 区域，单光子未分辨窗口的无量纲权重积分为
\begin{equation}
 {
 \mathcal R(\lambda_{\rm IR},\Delta E_{\rm det})
 =\frac{\alpha}{4\pi^2}
 \int_{\lambda_{\rm IR}/\hbar}^{\Delta E_{\rm det}/\hbar}
 \frac{\dd\omega}{\omega}
 \int\dd\Omega\,
 \left|
 \sum_i\sigma_i^{\rm io}Q_i
 \frac{\boldsymbol\beta_i^\perp}
 {1-\vb{n}\!\cdot\boldsymbol\beta_i}
 \right|^2.}
 \label{eq:soft-radiator-definition-dp10}
\end{equation}
频率积分显式含 $\ln(\Delta E_{\rm det}/\lambda_{\rm IR})$，因此“硬过程加一条真实软光子”的概率单独并没有 $\lambda_{\rm IR}\to0$ 极限。

同一微扰阶的虚软光子修正具有完全相同的外腿软因子核。把虚光子割到质量壳上会产生与\eqnrefs{eq:soft-radiator-definition-dp10} 相同的一光子相空间，而幺正性固定其概率流方向为相反号。红外奇异部分满足
\begin{equation}
 {
 2\Re\frac{\delta_{\rm virt}^{\rm soft}\mathcal M_0}{\mathcal M_0}
 =-\mathcal R(\lambda_{\rm IR},\Delta E_{\rm det})
 +C_{\rm IR\,finite},}
 \label{eq:virtual-real-ir-relation-dp10}
\end{equation}
其中 $C_{\rm IR\,finite}$ 在 $\lambda_{\rm IR}\to0$ 时有限。于是包含未分辨真实软辐射的包容概率满足
\begin{equation}
 {
 \lim_{\lambda_{\rm IR}\to0}
 P_{\rm incl}(\Delta E_{\rm det})<\infty.}
 \label{eq:bloch-nordsieck-cancellation-dp10}
\end{equation}
调节参数从物理预测中消失，但 $\Delta E_{\rm det}$ 保留下来，因为它属于实验测量定义。

\begin{proposition}[独立软发射权重的指数归一化]
固定一个硬散射事件，设红外截止下的单光子权重 $0\le\mathcal R<\infty$。
每个光子分别位于同一指定软窗口，积分区域为单光子区域的笛卡尔积，不施加总软能量切选。
假定软发射独立，$P_n^{\rm soft}$ 是该硬事件条件下归一化的光子数概率，则不可分辨光子的组合权重与归一化给出以下 Poisson 分布。
领先软独立发射模型的 $n$ 光子权重满足
\begin{equation}
 {
 P_n^{\rm soft}=\frac{\mathcal R^n}{n!}P_0^{\rm soft}.}
 \label{eq:soft-independent-n-emission-dp68}
\end{equation}
幺正性再固定零光子概率与全部 $n$ 光子概率的共同归一化：
\begin{equation}
 {
 P_0^{\rm soft}=\ee^{-\mathcal R},
 \qquad
 P_n^{\rm soft}
 =\ee^{-\mathcal R}\frac{\mathcal R^n}{n!}.}
 \label{eq:bloch-nordsieck-exponentiation-dp10}
\end{equation}
这就是 Bloch--Nordsieck 指数化：虚软修正降低严格无辐射通道的概率，真实软发射把同一概率重新分配到实验上退化的多光子扇区。
\end{proposition}
\begin{proof}
领先软极限中，不同软光子从同一个经典化外腿电流独立发射。若单个软窗口的积分为 $\mathcal R$，则 $n$ 个不可分辨光子的组合因子为 $1/n!$，所以

\begin{equation*}
 P_n^{\rm soft}=\frac{\mathcal R^n}{n!}P_0^{\rm soft}.
\end{equation*}
幺正性要求所有不可分辨光子数的概率和为一：
\begin{align*}
 1
 &=\sum_{n=0}^{\infty}P_n^{\rm soft}
 =P_0^{\rm soft}\sum_{n=0}^{\infty}\frac{\mathcal R^n}{n!},\\
 &=P_0^{\rm soft}\ee^{\mathcal R}.
\end{align*}
因此
\begin{equation*}
 P_0^{\rm soft}=\ee^{-\mathcal R},
 \qquad
 P_n^{\rm soft}=\ee^{-\mathcal R}\frac{\mathcal R^n}{n!},
\end{equation*}
即\eqnrefs{eq:bloch-nordsieck-exponentiation-dp10}。
该幂级数求和在给定独立软模型内严格成立，不等同于含任意探测器切选的完整 QED 包容概率。
\end{proof}

\begin{remark}[总软能量切选与独立性]
若实验要求 $\sum_{a=1}^n\hbar\omega_a<\Delta E_{\rm det}$，积分区域不再是单光子区域的笛卡尔积。
即使发射振幅仍取领先软因子，也须在多光子相空间内加入该共同切选，不能直接采用 $\mathcal R^n/n!$。
因此总能量切选本身的效应并非必然属于次领先软修正；逐光子窗口 Poisson 模型应与实际包容测量规则区分。
\end{remark}

有质量 QED 中，电子质量切断严格共线发散；当 $m_{\mathrm e} c^2$ 远小于硬能标时，共线区仍产生大的质量对数。若电子能量为 $E$、光子与电子夹角为 $\theta$，则
\begin{equation}
 {
 p\!\cdot k
 =\frac{E\hbar\omega}{c^2}
 \left(1-\beta_{\rm rel}\cos\theta\right).}
 \label{eq:soft-collinear-dot-exact-dp68}
\end{equation}
在 $\theta\ll1$ 且 $m_{\mathrm e} c^2/E\ll1$ 时，才进一步展开为
\begin{equation}
 {
 p\!\cdot k
 \simeq
 \frac{E\hbar\omega}{2c^2}
 \left(\theta^2+\frac{m_{\mathrm e}^2c^4}{E^2}\right).}
 \label{eq:soft-collinear-denominator-dp10}
\end{equation}
从而软极限与共线极限可以重叠。Kinoshita--Lee--Nauenberg 定理要求对给定实验分辨率下所有退化的初态和末态作适当求和或平均；Bloch--Nordsieck 对未分辨末态软光子的求和只是其中一个特例。

\begin{note}[带电渐近态与 LSZ]
严格无质量光子 QED 中，带电渐近态携带无限软光子云，局域裸电子场与孤立单粒子极点之间的朴素对应会被红外粒子结构改变。固定阶散射计算可先引入红外调节参数或使用适当缀饰态，在形成包容可观测量以后再取去调节极限。由此，LSZ 的孤立极点假设与严格红外极限的带电渐近态不发生逻辑冲突。
\end{note}

外部在壳极点产生普适软因子，光子相空间把 $1/\omega$ 振幅转化为红外对数，虚软程函核由幺正性给出相反的红外对数，包容求和消去调节参数，多软光子概率再由同一幺正性结构指数化。QCD 中对应的软与共线因子化保留这条动力学链，但电荷数被不对易的颜色生成元取代。

